% Planned 202.3--202.4 continuation; follows the current scheme/CQ foundations.
\section{From an example to a generalization: expose the extrapolation}
\label{sec:p18-example-sign-schemes}

An example can establish that a configuration is possible, illustrate a rule
already accepted, suggest a generalization, or support a judgment about another
case. These conclusions have different force. The 2017 Macagno--Walton--Reed
classification places argument from example under induction, while warning
that material relations and inference types cross-cut. That table label alone
does not supply an extrapolation warrant or a complete question list
\parencite[pp.~2495--2496, 2518--2519]{MacagnoWaltonReed2017Schemes}.
The following reconstruction makes the extrapolation warrant explicit;
its assessment prompts are editorial rather than a quotation of a
catalogue's question set.

\begin{hbdefinition}{def:p18-example-extrapolation}{The instance, scope, and needed warrant}
Fix a proposed reference class \(R\), condition \(F\), outcome \(G\), and
recorded case \(a\). The case premise says that \(a\) really is in \(R\)
and has both \(F\) and \(G\), under recorded circumstances. The proposed
conclusion says that other \(F\)-cases in \(R\) generally have \(G\), with
a specified qualifier. A needed warrant explains why the selected case or
sample is informative about those other cases: a defensible selection design,
stable common mechanism, relevant similarity, or separately supported default.
If that warrant is absent, the case supplies an observation and a candidate
hypothesis, not a universal rule. A warrant that simply repeats the desired
generalization has not established it.
\end{hbdefinition}

For assessment, check five different matters. First verify the reported case:
its provenance, measurements, and both predicates. Second test its relevance:
does it actually instantiate the proposed condition, outcome, and scope?
Third examine typicality through the selection process, not through resemblance
to a memorable story. Fourth fix strength: possible, usual, a numerical rate,
or exceptionless universal are not interchangeable. Fifth inspect special
circumstances that could account for the outcome without transferring it.
Each answer needs evidence. These prompts form an operational assessment;
they are not attributed as a source-specific numbered question list.

\subsection{A complete example argument, then its failed expansion}
\label{sec:p18-example-worked}

\begin{hbexample}{def:p18-example-decoder-case}{One successful conversion is not a release-wide guarantee}
All records in this case are hypothetical. An archive converts one image
\(a\) with Decoder V. It is an ordinary uncompressed image with profile P;
the input checksum, byte-level decoded pixels, dimensions, and profile handling
are independently checked. The claim ``Decoder V can correctly decode at least
one profile-P image'' is supported by that verified witness. It is an existence
claim, not a population-frequency claim.

The engineer instead concludes: ``Decoder V correctly decodes every image in
our collection.'' The case premise is true, but the conclusion changes the
reference class from uncompressed profile-P images to a collection containing
compressed, high-bit-depth and differently tagged images. There is no evidence
that the chosen file was selected independently of known successes. A decoder
that handles P but rejects the other formats is an explicit countercase:
the offered premise can be true while the universal conclusion is false.
The failure is in the extrapolation, not in the verified success.

Now stipulate a narrower assessment. The archive defines its population as
the profile-P records present at the frozen release date, assigns every record
a unique identifier, and draws thirty identifiers uniformly and independently
with replacement \emph{before} conversion results are known. All thirty
passes use the same declared success criterion. The checksum comparison
establishes each observed success. Duplicate draws are recorded, not
misrepresented as thirty distinct files. Nothing is inferred about other
formats or a later changed decoder.

Stipulate also that the frozen decoder is stateless and deterministic on
these records. Success is then a fixed property of each identifier.
Independent uniform draws give independent Bernoulli success indicators,
with a common probability equal to the fraction of identifiers that pass.
Uniform identifier selection alone would not establish this model for a
stateful decoder or a shared randomly failing conversion environment.

Case truth is supported by the checks. Relevance is fixed by P and the success
criterion. Typicality is supported by the prospective uniform design, not by
the engineer's preferred examples. Strength is a statistical test statement
under that design, not ``all files pass.'' Special circumstances such as a
training-only directory or hand-repaired inputs are excluded by the stipulated
sampling frame and untouched inputs. The calculation below supplies the
precise limited result. If the thirty files were selected from a folder of
previously successful demonstrations, its independence-of-outcome premise
would fail and that result would not apply.
\end{hbexample}

\begin{hbtheorem}{thm:p18-example-all-pass-bound}{proved=annex:p18-example-all-pass-bound}{The exact scope of an all-pass sample}
Suppose \(n\geq1\) trials are independent Bernoulli observations with the
same success probability \(p\). For \(0<\varepsilon<1\), under the
condition \(p\leq1-\varepsilon\), the probability that all \(n\) trials
pass is at most \((1-\varepsilon)^n\).
Thus a predeclared procedure that rejects that condition only on all-pass
has false-rejection probability at most this bound under those assumptions.
Observing all-pass neither deductively proves \(p>1-\varepsilon\) nor
assigns a posterior probability to that assertion without further premises.
\end{hbtheorem}

For the stipulated thirty-draw case, \(\varepsilon=0.1\) gives
\(0.9^{30}<0.043\). That is a bound on a specified procedure's error
under a sampling model. It is not a measured decoder accuracy, a posterior
probability, or a Walton critical-question score. Increasing a biased sample's
size does not repair its selection warrant.

\subsection{Classification and analogy do different work}
\label{sec:p18-example-classification-analogy}

Suppose the archive defines a ``profile-P record'' by a declared header and
encoding criterion. Establishing those features for a file classifies that
file. Under an accepted strict definition and verified features, the inference
can be deductive. But ``profile-P files are normally decoded correctly'' is
a performance generalization, not the definition. Classification does not
silently make that empirical warrant exceptionless.

A visually similar file may have a different compression tag. Inferring its
type from appearance is a defeasible classification step; using the metadata
criterion is another route. A default category can also have exceptions.
The current factual/value classification foundation is therefore retained:
an object meeting a definition, an object predicted to perform in a certain
way, and a performance judged acceptable need different premises.

For analogy, take a second archive using the same decoder build and profile.
Its curator proposes the same limited pilot protocol, not a universal success
guarantee. Stipulate that both archives seek a bounded feasibility check, have
authorized access to the records and resources for the protocol, and accept
this limited pilot as an appropriate way of obtaining that information.
Those practical commitments support the proposed action; they do not follow
from a similarity of image formats alone.
Walton's 2014 source-to-target presentation requires a described
source case, a plausibly justified source conclusion, and relevant similarity.
Its three questions examine differences weakening that similarity, justification
of the source conclusion, and a relevant countercase with a different warranted
outcome \parencite[author pp.~1--3]{Walton2014AnalogySchemes}.
Apply them here: a different decoder build or success criterion defeats the
claimed comparison; an unverified source pilot cannot support transfer; a
collection with untested profiles is a relevant countercase. With all those
factors matched in the stipulated records, the analogy supports repeating a
bounded pilot, not copying the source's achieved success rate as a known fact.
This preserves the existing factor/script variants rather than conflating
analogy with a sampled population estimate.

\section{A standalone sign inference: finding, context, and indication}
\label{sec:p18-sign-entry}

\begin{hbdefinition}{def:p18-sign-walton-reed-2002}{The actual author-version sign scheme}
Walton and Reed's 2002 author presentation reproduces a sign scheme with
three roles: a finding \(A\) is true in this situation; \(B\) is generally
indicated when \(A\) is true in this kind of situation; and \(B\) is
therefore accepted in this situation. The middle premise is a presumptive
warrant, subject to exceptions, not a material conditional guaranteed in every
case. The authors' false-track possibility explicitly defeats an unqualified
inference from an observed mark to its supposed origin
\parencite[author pp.~4--5]{WaltonReed2002DiagrammingAuthor}.
\end{hbdefinition}

That presentation gives the sign form and defeasibility explanation, but
does not display a sign-specific numbered question list. A sign's
observation, contextual applicability, strength of
indication, and rival origins still require examination; the next completed
source-specific question set belongs to the distinct abductive step.

\begin{hbexample}{def:p18-sign-pool-case}{A wet patch and a suspected pipe leak}
Consider a fictional archive room. A camera image shows a bright patch.
First, a worker touches and samples the surface and checks the image under
another illumination: the patch really is water, not a reflection. This
supports the finding \(A\). It does not yet identify its cause.
The worker argues that a leak \(B\) probably occurred because fresh pools
under this pipe generally indicate leaks when the room has not been cleaned
and no open container is present. Those context restrictions are part of
the warrant, not optional decorations.

In the successful stipulated record, the last cleaning was two days earlier,
the floor was dry afterward, the patch formed during the following hour,
no container was present, and a joint above it is actively dripping. The
observation is verified, the relevant context is established, and the direct
joint observation supplies a stronger, independent reason for a present leak.
There is defeasible support for investigating and containing that leak.
Blaming the installer or predicting damage to every stored object would
require further causal and practical premises.

In the failure record, a cleaner has just emptied a bucket at that spot.
The camera and water finding remain accurate, but the no-cleaning condition
fails and a concrete alternative origin is established. The sign inference
alone no longer supports the leak conclusion. It does not establish that
there is no leak: cleaning and a leak could coexist. A timed maintenance
light provides a different failure case: its lit state indicates that its
timer expired by convention, not that water exists or a pipe failed.
The operation of the indicator must be identified before a causal story
is supplied for its appearance.
\end{hbexample}

\section{From a sign to a diagnosis: Walton's distinct abductive scheme}
\label{sec:p18-sign-abductive-diagnosis}

Walton's 2001 scheme has three premises: \(F\) is a finding or presently
accepted fact set; \(E\) satisfactorily explains \(F\); and no alternative
offered so far explains it as satisfactorily. Its conclusion is that \(E\)
is a plausible hypothesis. ``Offered so far'' and the hypothesis qualifier
are essential. It is not a proof that the candidate list is exhaustive
\parencite[pp.~160--163]{Walton2001AbductivePresumptive}.
The complete four-question list of that presentation asks:

\begin{enumerate}
\item How satisfactory is \(E\) on its own, before comparing the available rivals?
\item How much better is it than the rivals actually available?
\item How far has the dialogue progressed, and how thorough is the inquiry's search?
\item Is further inquiry preferable to drawing a conclusion now?
\end{enumerate}

In the successful pool record, the leak explains a freshly appearing pool
and the drip without adding a new unobserved event. Cleaning fails on the
recorded timing; a spilled container fails on the checked absence and continuing
drip. The investigation has checked those concrete alternatives, not every
logically possible cause. A short temporary containment action can proceed
while a maintenance examination continues; a final responsibility attribution
should wait. This answers the first two questions with the actual mechanism
and comparisons, the third with the search performed, and the fourth with
the costs and reversibility of the provisional action.

If the only fact is an alarm, ``the pump failed'' may explain it but not be
better than a sensor fault, blocked channel, or high load. If a cheap pump-speed
check is available, the fourth question favors continuing inquiry. Later
verification that the pump is running and the sensor is stuck changes the
assessment: the original hypothesis is withdrawn, rather than protected by
inventing another premise. Walton's note~5 also warns that a predictive sign,
such as clouds forecasting rain, need not be an explanation of a past finding.
Thus these four abductive questions are not advertised as every sign scheme's
universal question set.

\section{A formal bridge for probabilistic signs, not an automatic cause}
\label{sec:p18-sign-probability}

\begin{hbdefinition}{def:p18-sign-probability-model}{Reference class and likelihoods}
In a declared probability model let \(H\) mean the hypothesized state and
\(S\) the observed sign. Put \(\pi=P(H)\), \(s=P(S\mid H)\), and
\(f=P(S\mid\neg H)\), with \(0<\pi<1\) and
\(s\pi+f(1-\pi)>0\). These values must concern the same reference
class, observation procedure, and context. A causal statement that H can
produce S does not supply \(f\), \(\pi\), or the reverse conditional.
\end{hbdefinition}

\begin{hbtheorem}{thm:p18-sign-bayes-bridge}{proved=annex:p18-sign-bayes-bridge}{A sign's posterior strength and alternatives}
Under the stated probability model,
\[
 P(H\mid S)=\frac{s\pi}{s\pi+f(1-\pi)}.
\]
For \(f>0\), posterior odds equal prior odds multiplied by \(s/f\).
For \(0<\tau<1\), the posterior is at least \(\tau\) exactly when
\((1-\tau)s\pi\geq\tau f(1-\pi)\). If \(s=f>0\), the sign leaves
the prior unchanged. A sign can therefore be strongly associated with H
yet give a weak posterior when alternative origins are frequent or H rare.
\end{hbtheorem}

In a hypothetical calibrated alert system, \(\pi=0.01,s=0.9,f=0.09\)
give \(P(H\mid S)=0.009/0.0981=10/109\), about \(9.2\%\), not
\(90\%\). The sign is ten times more likely under H, but the non-H
class is ninety-nine times larger. With otherwise identical likelihoods
and \(\pi=0.5\), the posterior is \(10/11\). These two models are
explicit countercases to using sensitivity alone as diagnostic certainty.
No number here is an empirical estimate for real archives or equipment.

Two alerts copied from one sensor are not independent corroboration.
Multiplying likelihood ratios requires a separately justified conditional
independence model; repetition of the same report does not supply it.
Nor does a posterior about H establish that intervening on H would change
S. The current ten-question correlation-to-cause entry remains the route
for temporal direction, common causes, mediation, measurement changes,
mechanism, and surrogate endpoints. The present bridge supplies an exact
conditional calculation, not those missing causal warrants.

\subsection{A fully specified correlation with no intervention effect}
\label{sec:p18-sign-common-cause-countercase}

Walton's causal discussion distinguishes evidence of association from its
direction and a possible common cause
\parencite[pp.~185--187]{Walton2016ArgumentEvaluation}.
The following fictional process gives a complete countercase rather than
merely mentioning confounding as an unperformed question.

\begin{hbdefinition}{def:p18-sign-confounding-process}{The independent-input device model}
Let \(C\) have values high and low, each with probability one half.
Let \(U_X,U_Y\) be independent uniform draws from \(\{1,\ldots,10\}\),
also independent of C. A monitoring retrofit X is installed exactly when
\(C=\mathrm{high},U_X\leq9\) or \(C=\mathrm{low},U_X\leq1\).
A device failure Y occurs exactly when \(C=\mathrm{high},U_Y\leq6\)
or \(C=\mathrm{low},U_Y\leq1\). An intervention setting X replaces
only the installation equation; it leaves the distributions of C and the
inputs and the failure equation unchanged.
\end{hbdefinition}

\begin{hbtheorem}{thm:p18-sign-confounding-countermodel}{proved=annex:p18-sign-confounding-countermodel}{Positive association does not identify the effect}
In this specified process, \(P(Y\mid X)=0.55\) and
\(P(Y\mid\neg X)=0.15\), but setting X present or absent both
produce failure probability \(0.35\). Thus a genuine positive observational
association is compatible with no intervention effect of X in this model.
\end{hbtheorem}

The observation argument is ``retrofitted devices fail more often, so the
retrofit causes failure.'' Its comparison is real \emph{within the stipulated
model}; it is not fabricated away. But high load makes both installation and
failure more frequent. Within each load stratum the failure rates are equal
for installed and uninstalled devices: \(0.6\) in high load and \(0.1\)
in low load. Installation was targeted to the riskier devices.

The descriptive marginal rates are supported. The causal conclusion fails
in the fully supplied alternative mechanism, which the annex calculates.
A remedy needs evidence identifying the intervention mechanism, not another
copy of the marginal rate. The failed claim also does not prove that every
possible retrofit has zero effect: a different installation/failure mechanism
could fit other records. The countermodel refutes the unrestricted
correlation-to-cause inference, while preserving the need to investigate the
actual mechanism in a real case.

\begin{proofboundary}{def:p18-example-sign-original-boundary}{The warrants of the three different inferences}
The sampling bound requires its stated independent, equal-probability
trial model. The sign calculation requires the same reference class and
the declared likelihoods and prior. Neither supplies a causal mechanism,
an exhaustive set of explanations, or a universally valid presumptive rule.
The 2002 sign form, 2001 abductive questions and 2014 analogy questions
belong to distinct source presentations; their questions and conclusion
qualifiers must not be combined into an unattributed universal scheme.
The example-assessment prompts here remain an editorial reconstruction.
\end{proofboundary}
